\section{本章小结}
\label{sec:ot-summary}

本章从一个问题出发：两个边缘分布固定，怎样把一个搬成另一个最省。可行方案是全部耦合，它们构成一个凸集；在有限情形是运输多面体，在一般情形靠固定边缘带来的胎紧性获得紧性。成本非负下半连续时最优耦合存在；去掉下半连续，极小化序列可能在极限处跳高。存在的是最优耦合，不是最优映射：映射要求质量不可拆分，可能连可行解都没有，或者下确界取不到。

最优性的证据来自对偶。满足$\varphi(x)+\psi(y)\leq c(x,y)$的一对势给出一切耦合成本的公共下界，原对偶差等于耦合下松弛量的积分，所以质量只落在零松弛处的耦合就是最优的。有限情形的强对偶由线性规划得到，紧致空间上的连续成本可以由有限问题逼近得到；原问题达到、无对偶间隙与最优势存在是三个不同的结论，需要分别核对条件。

具体结构决定怎样求解。表~\ref{tab:ot-summary}汇总本章求解的几类问题。

\begin{table}[htbp]
 \centering
 \caption{本章求解的运输问题。高斯两行为平方成本，实数轴一行要求$p\geq1$与有限$p$阶矩；“映射”一列指最优方案是否由运输映射实现。}
 \label{tab:ot-summary}
 \begin{tabular}{lll}
  \toprule
  问题 & 最优方案 & 映射\\
  \midrule
  实数轴，$\lvert x-y\rvert^{p}$ & 分位数耦合 & 源无原子时是\\
  有限支持 & 线性规划的顶点 & 一般不是\\
  高斯，$\bm{\Sigma}_1$正定 & 对称正定仿射映射 & 是\\
  高斯，$\bm{\Sigma}_1$奇异 & 按支撑子空间分析 & 可能不是\\
  熵正则化 & 行列缩放式~\eqref{eq:ot-scaling-form} & 不是\\
  \bottomrule
 \end{tabular}
\end{table}

一维问题证明了第~\ref{chap:distribution-comparison}章提到的关系：$W_1$等于两条分布函数之间的面积，$W_p^{p}$等于分位数函数之差的$p$次方积分。全章例子中，逐格搬运同时被对偶势、分位数耦合证实为平方成本下唯一的最优方案，成本$0.6$。熵正则化把运输问题变成严格凸的光滑问题，解具有行列缩放形式，Sinkhorn迭代交替满足行和与列和；正则化偏差不超过$\varepsilon$乘以最优方案与独立方案之间的KL散度，迭代误差由修正可行方案后的原对偶差控制，两者与离散化误差要分别报告。

最后两节用最优耦合构造新的分布。位移插值把每份质量沿直线匀速推向目的地，是$W_2$下的常速测地线；混合路径则只调整权重，在$W_2$下甚至不是有限长的路径。Wasserstein重心推广了点的加权平均，在有限二阶矩下存在；一维时就是分位数函数的加权平均，高维时唯一性需要某个正权重分布绝对连续。最优耦合不唯一时，插值与重心都不唯一。

本章始终在成本、矩条件、映射可行性与数值误差的范围内下结论：有限$p$阶矩保证成本有限，下半连续保证最优耦合存在，严格凸成本与绝对连续的源分布才带来唯一性与映射，正则化与迭代给出的是带有可计算误差界的近似。$W_p$作为距离的性质及其与弱收敛的关系，分别见第~\ref{chap:distribution-comparison}章与第~\ref{chap:weak-convergence-of-distributions}章。
